\documentclass[11pt,a4paper]{article}
\usepackage[utf8]{inputenc}
\usepackage[T1]{fontenc}
\usepackage{lmodern}
\usepackage[a4paper,margin=1in]{geometry}
\usepackage{amsmath,amssymb,amsfonts,amsthm}
\usepackage{graphicx}
\usepackage{tikz}
\usetikzlibrary{arrows.meta, positioning, decorations.pathmorphing, shapes.geometric, calc}
\usepackage{tcolorbox}
\tcbuselibrary{skins,breakable}
\usepackage{enumitem}
\usepackage{booktabs}
\usepackage{tabularx}
\usepackage{array}
\usepackage{xcolor}
\usepackage{titlesec}
\usepackage{physics}
\usepackage{hyperref}
\hypersetup{
    colorlinks=true,
    linkcolor=blue!60!black,
    citecolor=red!60!black,
    urlcolor=blue!70!black
}
\usepackage{cleveref}

% ---------------------------------------------------------------------
% Macros
% ---------------------------------------------------------------------
\newcommand{\ee}{\mathrm{e}}
\newcommand{\ii}{\mathrm{i}}
\newcommand{\kB}{k_{\mathrm{B}}}
\renewcommand{\dd}{\mathrm{d}}
\newcommand{\calL}{\mathcal{L}}
\newcommand{\calH}{\mathcal{H}}
\newcommand{\calD}{\mathcal{D}}
\newcommand{\calC}{\mathcal{C}}
\newcommand{\calM}{\mathcal{M}}
\newcommand{\bbR}{\mathbb{R}}
\newcommand{\bbC}{\mathbb{C}}
\newcommand{\bbZ}{\mathbb{Z}}
\newcommand{\bbN}{\mathbb{N}}
\newcommand{\Cz}{C_z}
\newcommand{\Bz}{B_z}
\newcommand{\gz}{g_z}

\theoremstyle{definition}
\newtheorem{definition}{Definition}[section]
\newtheorem{example}{Example}[section]
\theoremstyle{remark}
\newtheorem{remark}{Remark}[section]

% Callout boxes for the three recurring pedagogical questions:
%  WHERE / WHY / WHEN an approximation is made, and HOW to implement it.
\newtcolorbox{approxbox}[1][]{
  colback=orange!6, colframe=orange!55!black, boxrule=0.6pt,
  breakable, left=6pt,right=6pt,top=4pt,bottom=4pt,
  title={\bfseries Approximation --- #1}}
\newtcolorbox{implbox}[1][]{
  colback=blue!5, colframe=blue!45!black, boxrule=0.6pt,
  breakable, left=6pt,right=6pt,top=4pt,bottom=4pt,
  title={\bfseries How to implement --- #1}}
\newtcolorbox{keybox}{
  colback=green!6, colframe=green!45!black, boxrule=0.8pt,
  breakable, left=6pt,right=6pt,top=4pt,bottom=4pt}

\title{\bfseries A Complete, From--Scratch Derivation of the\\
Thermodynamic Neuron\\[4pt]
\large From a single qubit to the autonomous NOT gate and the general
linearly--separable neuron}
\author{Sparsho Chakraborty\\[2pt]
\small (pedagogical companion to Lipka--Bartosik, Perarnau--Llobet \&
Brunner, \emph{Sci.\ Adv.} \textbf{10}, eadm8792, 2024)}
\date{\today}

\begin{document}
\maketitle

\begin{abstract}
\noindent
This note rebuilds the ``thermodynamic neuron'' of
Ref.~\cite{LipkaBartosik2024} from the ground up, at the level of detail
a student needs in order to reproduce every equation by hand and in
code. We start from a single qubit in a heat bath and derive, step by
step: (i) the thermal (Gibbs) state and the Fermi--Dirac/sigmoid
occupation; (ii) the open--system master equation, showing explicitly
\emph{where the bath Hamiltonian goes} and \emph{why there is no
surviving qubit--bath interaction Hamiltonian in the final equations}
--- the single most common point of confusion; (iii) the reset
(collisional) thermalization model and its heat current; (iv) the
three--qubit machine, the \emph{virtual qubit}, and the \emph{virtual
temperature}; (v) the full interaction Hamiltonian of the collector and
why it must be a three--body, energy--conserving swap; (vi) why the
collector alone is only a linear amplifier and hence \emph{why the
modulator is needed}; (vii) the two--time--scale dynamics, the
steady--state transfer characteristic, and the emergence of the sigmoid;
(viii) logic encoding, readout error, and the error--dissipation
trade--off; and (ix) the generalization to any linearly--separable
Boolean function and its exact correspondence with the perceptron. Every
approximation is called out in an orange box (\emph{where, why, when}),
and every result is accompanied by a blue box explaining \emph{how to
implement it} in QuTiP. Equation numbers of the source paper are quoted
throughout as ``[paper Eq.~(n)]''.
\end{abstract}

\tableofcontents
\newpage

% =====================================================================
\part*{Reading guide}
\addcontentsline{toc}{part}{Reading guide}

The derivation is organized as a ladder. Each rung is self--contained and
ends with the object the next rung needs.

\begin{center}
\begin{tabular}{@{}rl@{}}
\toprule
Part & What is built \\
\midrule
I   & One qubit $\to$ open--system master equation $\to$ reset model
      $\to$ multi--qubit machine. \\
II  & Two qubits $\to$ virtual qubit $\to$ virtual temperature $\to$
      resonant three--body interaction. \\
III & Collector $+$ modulator $\to$ dynamics $\to$ transfer curve
      $\to$ sigmoid $\to$ NOT gate $\to$ error/dissipation. \\
IV  & $n$--input neuron $\to$ perceptron equivalence $\to$ design
      algorithm $\to$ NOR, MAJORITY, XOR networks. \\
\bottomrule
\end{tabular}
\end{center}

\noindent
Three questions are answered repeatedly and are flagged with coloured
boxes:
\begin{itemize}[leftmargin=*]
\item \textbf{\textcolor{orange!70!black}{Approximation boxes}} state
\emph{which} approximation is made, \emph{why} it is legitimate, and
\emph{when} (in which parameter regime) it holds.
\item \textbf{\textcolor{blue!60!black}{Implementation boxes}} translate
each result into concrete QuTiP objects.
\end{itemize}

% #####################################################################
\part{Physical foundations: from one qubit to a thermal machine}
\label{part:foundations}
% #####################################################################

\section{The single qubit}
\label{sec:single-qubit}

\subsection{Hilbert space, states, Hamiltonian}
A qubit is a two--level quantum system. Its Hilbert space is
$\calH = \bbC^2$, with an orthonormal basis given by the ground state
$\ket{0}$ and the excited state $\ket{1}$,
\begin{equation}
\ket{0} = \begin{pmatrix}1\\0\end{pmatrix},\qquad
\ket{1} = \begin{pmatrix}0\\1\end{pmatrix},\qquad
\braket{i}{j} = \delta_{ij}.
\end{equation}
A general (possibly mixed) state is a density operator $\rho$: a
positive semidefinite, Hermitian, unit--trace operator,
\begin{equation}
\rho = \rho^\dagger \succeq 0,\qquad \Tr\rho = 1 .
\end{equation}
The free Hamiltonian assigns energies $E_0$ and $E_1>E_0$ to the two
levels,
\begin{equation}
H = E_0\ket{0}\bra{0} + E_1\ket{1}\bra{1}.
\label{eq:H-two-energies}
\end{equation}
Only \emph{energy differences} are physical, so we are free to set the
zero of energy at the ground state, $E_0 = 0$. Then
\begin{equation}
\boxed{\,H = \epsilon\,\ket{1}\bra{1} = \epsilon\, n,\qquad
\epsilon := E_1 - E_0,\quad n := \ket{1}\bra{1}.\,}
\label{eq:H-single}
\end{equation}
The single number $\epsilon$ --- the \emph{energy gap} --- completely
specifies the qubit. This is exactly the Hamiltonian
$H = E_0\ket{0}\bra{0}+E_1\ket{1}\bra{1}$ of [paper, Sec.~II\,A\,2] with
$E_0=0$.

\begin{remark}
It is often convenient to use Pauli operators. With
$\sigma_z = \ket{0}\bra{0}-\ket{1}\bra{1}$ one has
$n = \tfrac12(\mathbb{1}-\sigma_z)$, so
$H = \tfrac{\epsilon}{2}(\mathbb{1}-\sigma_z)$, which equals
$-\tfrac{\epsilon}{2}\sigma_z$ up to the physically irrelevant constant
$\tfrac{\epsilon}{2}\mathbb{1}$. We keep the form $\epsilon\, n$ because
the ``excited--state population'' $n$ is the quantity we will read out.
The raising/lowering operators are
$\sigma_+ = \ket{1}\bra{0}$, $\sigma_- = \ket{0}\bra{1}$, with
$\sigma_x=\sigma_++\sigma_-$.
\end{remark}

\subsection{The thermal (Gibbs) state and the partition function}
Suppose the qubit is left in contact with a large heat bath at inverse
temperature $\beta = 1/(\kB T)$ for a long time. Statistical mechanics
tells us it settles into the canonical (Gibbs) state
\begin{equation}
\tau(\beta) := \frac{\ee^{-\beta H}}{Z},\qquad
Z := \Tr\,\ee^{-\beta H} = 1 + \ee^{-\beta\epsilon}.
\label{eq:gibbs}
\end{equation}
Because $H$ is diagonal, so is $\tau(\beta)$:
\begin{equation}
\tau(\beta) = \frac{1}{1+\ee^{-\beta\epsilon}}\ket{0}\bra{0}
            + \frac{\ee^{-\beta\epsilon}}{1+\ee^{-\beta\epsilon}}
              \ket{1}\bra{1}.
\label{eq:gibbs-explicit}
\end{equation}
This is [paper Eq.~(8)] for a single qubit.

\subsection{Excited population = Fermi--Dirac = sigmoid}
The probability that the qubit is found excited is the $\ket{1}\bra{1}$
matrix element of $\tau(\beta)$,
\begin{equation}
g(\beta\epsilon) := \bra{1}\tau(\beta)\ket{1}
= \frac{\ee^{-\beta\epsilon}}{1+\ee^{-\beta\epsilon}}
= \frac{1}{1+\ee^{\beta\epsilon}} .
\label{eq:fermi}
\end{equation}
The last equality (multiply numerator and denominator by
$\ee^{+\beta\epsilon}$) is [paper Eq.~(2)]. Three facts about $g$ will be
used constantly:
\begin{enumerate}[label=(\alph*),leftmargin=*]
\item \textbf{It is the Fermi--Dirac function}, monotonically decreasing
in $\beta\epsilon$, with $g(-\infty)=1$, $g(0)=\tfrac12$,
$g(+\infty)=0$.
\item \textbf{It is the logistic sigmoid} $f(x) = (1+\ee^{x})^{-1}$
evaluated at $x=\beta\epsilon$. This single identity is the seed of the
entire neuron--perceptron analogy: a thermal population \emph{is} a
sigmoid of an (inverse) temperature.
\item \textbf{Detailed--balance ratio.} From \eqref{eq:gibbs-explicit},
\begin{equation}
\frac{\bra{1}\tau(\beta)\ket{1}}{\bra{0}\tau(\beta)\ket{0}}
= \ee^{-\beta\epsilon}.
\label{eq:db-ratio}
\end{equation}
The ratio of excited to ground population is the Boltzmann factor. This
is the microscopic fingerprint of temperature and is the property every
thermalization model must reproduce.
\end{enumerate}

\begin{implbox}[the thermal state]
In QuTiP: \texttt{H = eps * create(2)*destroy(2)} (i.e.\ $\epsilon\,n$),
and the Gibbs state is \texttt{tau = (-beta*H).expm(); tau/=tau.tr()}.
The excited population is \texttt{expect(num(2), tau)}, which returns
$g(\beta\epsilon)$ of \eqref{eq:fermi}. Equivalently
\texttt{g = 1/(1+np.exp(beta*eps))}.
\end{implbox}

% ---------------------------------------------------------------------
\section{Open quantum systems: where the bath Hamiltonian goes}
\label{sec:open}

This section answers the question your reader raised: \emph{if the qubit
is coupled to a bath, where is the qubit--bath interaction Hamiltonian in
the equations, and why does it disappear?} The short answer: it is still
there microscopically, but after we (i) enlarge the description to
system$+$bath, (ii) assume weak coupling and short bath memory, and
(iii) trace the bath out, the interaction Hamiltonian is \emph{converted}
into the dissipator $\calL[\rho]$ --- a term acting on the system alone.
We now show this explicitly.

\subsection{Why a closed qubit never thermalizes}
A closed quantum system evolves unitarily, $\rho(t) = U(t)\rho(0)
U^\dagger(t)$ with $U=\ee^{-\ii Ht}$. For the single qubit
\eqref{eq:H-single}, $[\rho,H]$ generates only phases: the populations
$\bra{0}\rho\ket{0}$, $\bra{1}\rho\ket{1}$ are constant in time.
A closed qubit therefore \emph{cannot} relax to the Gibbs state
\eqref{eq:gibbs}; it just keeps whatever populations it started with.
Thermalization, dissipation, and irreversibility are impossible without
an environment. This is the physical reason a bath is not optional: it
is the only object in the model that can change populations and drive
the system to a steady state.

\subsection{The microscopic system$+$bath$+$interaction Hamiltonian}
We restore the environment explicitly. Model the bath as a large
collection of harmonic oscillators (the standard Caldeira--Leggett /
quantum--optics reservoir),
\begin{equation}
H_B = \sum_k \omega_k\, a_k^\dagger a_k ,
\label{eq:HB}
\end{equation}
with bosonic modes $[a_k,a_{k'}^\dagger]=\delta_{kk'}$. The qubit couples
to the bath \emph{bilinearly and transversally},
\begin{equation}
H_{SB} = \sigma_x \otimes B, \qquad
B := \sum_k \lambda_k\,(a_k + a_k^\dagger),
\label{eq:HSB}
\end{equation}
where $B$ is a bath ``force'' operator and $\lambda_k$ are coupling
constants. The \emph{total} Hamiltonian on $\calH_S\otimes\calH_B$ is
\begin{equation}
\boxed{\,H_{\mathrm{tot}} = \underbrace{H_S}_{\epsilon\, n}
\;+\; \underbrace{H_B}_{\sum_k\omega_k a_k^\dagger a_k}
\;+\; \underbrace{H_{SB}}_{\sigma_x\otimes B}\, .}
\label{eq:Htot-open}
\end{equation}
Two modelling choices deserve a word, because they are exactly the
places where physics enters.

\begin{approxbox}[bilinear coupling, \eqref{eq:HSB}]
\textbf{Where:} the system--bath term is taken linear in the bath
operators, $H_{SB}=\sigma_x\otimes\sum_k\lambda_k(a_k+a_k^\dagger)$.
\textbf{Why:} this is the leading term in an expansion of a generic
coupling in bath displacements; it is \emph{exact} for a Gaussian
(harmonic) environment. \textbf{When:} valid whenever the environment
responds linearly (weak, non--saturating coupling), which is the regime
of local detailed balance we ultimately need.
\end{approxbox}

\begin{remark}[why $\sigma_x$ and not $\sigma_z$]
The qubit operator in \eqref{eq:HSB} is \emph{transverse}
($\sigma_x=\sigma_++\sigma_-$). It must be, because to thermalize the
bath has to be able to \emph{flip} the qubit between $\ket{0}$ and
$\ket{1}$, exchanging the energy quantum $\epsilon$. A longitudinal
coupling $\sigma_z\otimes B$ commutes with $H_S$ and only randomizes the
phase (pure dephasing): it can never change populations and thus never
produces a thermal state. Keep this contrast in mind --- when we later
add a Floquet drive we will deliberately make it \emph{longitudinal},
precisely so that it does \emph{not} pump populations.
\end{remark}

\subsection{Tracing out the bath: the Born--Markov--secular derivation}
\label{sec:BMS}
We want an equation for the reduced system state
$\rho(t) := \Tr_B\,\rho_{\mathrm{tot}}(t)$ alone. Move to the interaction
picture with respect to $H_0 := H_S+H_B$ (tilde denotes interaction
picture),
\begin{equation}
\tilde H_{SB}(t) = \ee^{\ii H_0 t} H_{SB} \ee^{-\ii H_0 t}
= \tilde\sigma_x(t)\otimes \tilde B(t).
\end{equation}
The exact (von Neumann) equation
$\dot{\tilde\rho}_{\mathrm{tot}}=-\ii[\tilde H_{SB}(t),
\tilde\rho_{\mathrm{tot}}]$ is integrated once and resubstituted, giving
the exact integro--differential equation
\begin{equation}
\dot{\tilde\rho}_{\mathrm{tot}}(t)
= -\ii[\tilde H_{SB}(t),\rho_{\mathrm{tot}}(0)]
  - \int_0^t \dd s\,
  [\tilde H_{SB}(t),[\tilde H_{SB}(s),\tilde\rho_{\mathrm{tot}}(s)]].
\label{eq:exact-integro}
\end{equation}
Tracing over the bath and applying three successive, clearly labelled
approximations turns this into a time--local master equation.

\begin{approxbox}[(A1) Born / weak coupling]
\textbf{Where:} we assume the state stays factorized,
$\rho_{\mathrm{tot}}(s)\approx \rho(s)\otimes\rho_B$, with $\rho_B$ the
fixed thermal state of the bath. \textbf{Why:} the coupling
$\lambda_k$ is weak, so system--bath correlations built up are
$O(\lambda^2)$ and negligible in the leading dissipator. \textbf{When:}
$\lambda$ (equivalently the induced rate $\gamma$) much smaller than the
system gap $\epsilon$ and the bath bandwidth. Choosing $\rho_B$ thermal
and $\Tr_B[B\rho_B]=0$ kills the first term of
\eqref{eq:exact-integro}.
\end{approxbox}

\begin{approxbox}[(A2) Markov / short bath memory]
\textbf{Where:} replace $\tilde\rho(s)\to\tilde\rho(t)$ in the integral
and send the lower limit to $-\infty$. \textbf{Why:} the bath
correlation function
$C(\tau)=\Tr_B[\tilde B(\tau)B\,\rho_B]$ decays on a microscopic time
$\tau_c$ much shorter than the system's relaxation time, so the system
has no memory of its past on the scale that matters. \textbf{When:}
$\tau_c \ll 1/\gamma$. This is the step that makes the equation
\emph{time--local} (memoryless).
\end{approxbox}

After (A1)--(A2), tracing out the bath yields the Redfield equation. Its
rates are the one--sided Fourier transforms of the bath correlation
function evaluated at the qubit's Bohr frequency $\pm\epsilon$:
\begin{equation}
\gamma(\omega) := \int_{-\infty}^{\infty}\dd\tau\,
\ee^{\ii\omega\tau}\,C(\tau)
= 2\pi\, J(\omega)\,[1+n_{\mathrm{B}}(\omega)],
\label{eq:rate-fourier}
\end{equation}
where $J(\omega)=\sum_k\lambda_k^2\delta(\omega-\omega_k)$ is the
\emph{spectral density} and $n_{\mathrm B}$ the Bose occupation. The
last, purely mathematical, step secures complete positivity.

\begin{approxbox}[(A3) secular / rotating--wave]
\textbf{Where:} drop terms rotating at $\ee^{\pm 2\ii\epsilon t}$ in the
interaction--picture dissipator. \textbf{Why:} they oscillate fast
compared to the slow relaxation $\gamma$ and average to zero over the
relaxation time. \textbf{When:} $\gamma\ll\epsilon$ (well--resolved
gap). This step guarantees the generator is of Lindblad (GKSL) form and
hence completely positive.
\end{approxbox}

The result, back in the Schr\"odinger picture, is the GKSL master
equation
\begin{equation}
\boxed{\;
\dot\rho = -\ii[H_S,\rho]
+ \gamma_\downarrow\Big(\sigma_-\rho\sigma_+
   -\tfrac12\{\sigma_+\sigma_-,\rho\}\Big)
+ \gamma_\uparrow\Big(\sigma_+\rho\sigma_-
   -\tfrac12\{\sigma_-\sigma_+,\rho\}\Big).
\;}
\label{eq:GKSL-qubit}
\end{equation}
This is the promised structure $\dot\rho = -\ii[H,\rho]+\calL[\rho]$ of
[paper Eq.~(1)]. \textbf{Here is the answer to the reader's question:}
the interaction Hamiltonian $H_{SB}=\sigma_x\otimes B$ is \emph{no longer
present as a Hamiltonian}. It has been converted, by the elimination of
the bath, into the two dissipative channels of \eqref{eq:GKSL-qubit}:
\begin{itemize}[leftmargin=*]
\item $\gamma_\downarrow$ multiplies $\sigma_-$ (emission: the qubit
gives a quantum $\epsilon$ to the bath);
\item $\gamma_\uparrow$ multiplies $\sigma_+$ (absorption: the qubit
takes a quantum $\epsilon$ from the bath).
\end{itemize}
The bath Hamiltonian $H_B$ has disappeared entirely; the coupling
$H_{SB}$ survives only through the two numbers
$\gamma_\downarrow,\gamma_\uparrow$. \emph{That} is why the working
equations of the paper contain a free Hamiltonian $H_S$ and a dissipator
$\calL$, but no qubit--bath interaction Hamiltonian: the latter was spent
to build $\calL$.

\subsection{Detailed balance fixes the ratio of the rates}
The two rates are not independent. The Kubo--Martin--Schwinger (KMS)
condition of a thermal bath forces
$\gamma(-\omega)=\ee^{-\beta\omega}\gamma(\omega)$; applied to
$\omega=\epsilon$,
\begin{equation}
\frac{\gamma_\uparrow}{\gamma_\downarrow} = \ee^{-\beta\epsilon}.
\label{eq:db-rates}
\end{equation}
This is \emph{local detailed balance}. Its whole purpose is to guarantee
that the steady state of \eqref{eq:GKSL-qubit} reproduces the Boltzmann
ratio \eqref{eq:db-ratio}. Indeed, writing $p:=\bra{1}\rho\ket{1}$, the
populations obey the classical rate equation
\begin{equation}
\dot p = \gamma_\uparrow(1-p) - \gamma_\downarrow\, p ,
\label{eq:pop-rate}
\end{equation}
whose fixed point $\dot p=0$ is
\begin{equation}
p_{\mathrm{ss}}
= \frac{\gamma_\uparrow}{\gamma_\uparrow+\gamma_\downarrow}
= \frac{\ee^{-\beta\epsilon}}{1+\ee^{-\beta\epsilon}}
= g(\beta\epsilon).
\end{equation}
The qubit relaxes to exactly the thermal excited population
\eqref{eq:fermi}, as it must.

\begin{keybox}
\textbf{Summary of Sec.~\ref{sec:open}.} Coupling a qubit to a bath
means writing $H_{\mathrm{tot}}=H_S+H_B+H_{SB}$; eliminating the bath
under weak coupling (A1), short memory (A2) and secular averaging (A3)
turns $H_{SB}$ into a dissipator $\calL[\rho]$ acting on the system
alone. The bath's only fingerprints in the final equation are the two
rates, whose ratio is pinned by detailed balance
$\gamma_\uparrow/\gamma_\downarrow=\ee^{-\beta\epsilon}$ so that the
steady state is the Gibbs state. This is why the paper never writes an
explicit qubit--bath interaction Hamiltonian.
\end{keybox}

\subsection{The reset model: a minimal thermalization dissipator}
\label{sec:reset}
Equation \eqref{eq:GKSL-qubit} is more detail than the logic of the gate
needs. Since only the \emph{populations} matter (the steady states are
diagonal), the paper adopts the simplest dissipator with the correct
fixed point: the \emph{reset model} [paper, Sec.~II\,A\,3],
\begin{equation}
\boxed{\;
\calL^{(k)}[\rho] = \gamma_k\Big(\Tr_k[\rho]\otimes\tau(\beta_k)
- \rho\Big).
\;}
\label{eq:reset}
\end{equation}
Its interpretation is collisional: in each infinitesimal interval, qubit
$k$ has probability $\gamma_k\,\dd t$ of being ``reset'' --- swapped for a
fresh qubit drawn from the bath's thermal state $\tau(\beta_k)$ ---
and probability $1-\gamma_k\,\dd t$ of being left alone. $\Tr_k[\rho]$
is the partial trace over qubit $k$ (the rest of the machine), so
$\Tr_k[\rho]\otimes\tau(\beta_k)$ is the state \emph{after} a full
thermalization event on qubit $k$.

\begin{approxbox}[(A4) reset thermalization, \eqref{eq:reset}]
\textbf{Where:} we replace the microscopically derived dissipator
\eqref{eq:GKSL-qubit} by the reset form \eqref{eq:reset}.
\textbf{Why:} both share the same fixed point (the local Gibbs state,
$\calL^{(k)}[\tau(\beta_k)]=0$) and the same detailed balance; the reset
model just fixes the relaxation to be a single rate $\gamma_k$ for all
observables. \textbf{When:} whenever only steady--state populations and
their thermodynamics are needed --- which is the whole computation.
Changing the thermalization model changes the exact nonlinear shape of
the transfer curve but not the machine's ability to invert
temperatures (this is stated explicitly around [paper Eq.~(23)]).
\end{approxbox}

For a single qubit, \eqref{eq:reset} reduces to
$\calL[\rho]=\gamma(\tau(\beta)-\rho)$, whose population equation is
$\dot p=\gamma[g(\beta\epsilon)-p]$ --- the same fixed point as before,
with relaxation rate $\gamma$.

\subsection{Heat current}
The heat delivered from a qubit to its bath is the energy the dissipator
removes from the system,
\begin{equation}
j_k := \Tr\!\big[H\,\calL^{(k)}[\rho]\big].
\label{eq:current-def}
\end{equation}
This is [paper Eq.~(5)]. For the reset model \eqref{eq:reset} with a
single qubit of gap $\epsilon_k$ and excited population $p_k$,
\begin{align}
j_k
&= \Tr\!\big[\epsilon_k n\,\gamma_k(\tau(\beta_k)-\rho)\big]
= \gamma_k\epsilon_k\big(\bra{1}\tau(\beta_k)\ket{1} - p_k\big)
\nonumber\\
&= \boxed{\,\gamma_k\epsilon_k\big[g(\beta_k\epsilon_k) - p_k\big].\,}
\label{eq:current-reset}
\end{align}
This is [paper Eq.~(6)]. The sign convention: $j_k>0$ when the bath
pumps energy into the qubit ($g(\beta_k\epsilon_k)>p_k$, the qubit is
colder than its bath). A qubit coupled to \emph{two} baths simply gets
the sum of two such dissipators, and the total current is the sum ---
a fact we will use for the modulator.

\begin{implbox}[the reset dissipator and current]
In QuTiP the reset dissipator is most easily built as a rate times the
difference of superoperators, or, for populations only, integrated as
the classical rate equation \eqref{eq:pop-rate}. For a full Lindblad
simulation one instead uses collapse operators
$\{\sqrt{\gamma_\downarrow}\,\sigma_-,\sqrt{\gamma_\uparrow}\,\sigma_+\}$
with $\gamma_\downarrow=\gamma(1-g)$, $\gamma_\uparrow=\gamma g$,
$g=g(\beta\epsilon)$; \texttt{mesolve} then reproduces
\eqref{eq:current-reset} as
\texttt{expect(H, rho)}'s rate of change due to that bath.
\end{implbox}

% ---------------------------------------------------------------------
\section{Multi--qubit thermal machines}
\label{sec:machine}

\subsection{Free and interaction Hamiltonian}
A machine is several qubits, labelled $k$, each with gap $\epsilon_k$.
The free Hamiltonian is the sum
\begin{equation}
H_0 = \sum_k \epsilon_k\,\ket{1}\bra{1}_k ,
\label{eq:H0-machine}
\end{equation}
where $\ket{1}\bra{1}_k$ acts as identity on all other qubits. The
qubits are coupled by a \emph{time--independent, energy--preserving}
interaction $H_{\mathrm{int}}$, meaning
\begin{equation}
[H_{\mathrm{int}}, H_0] = 0 .
\label{eq:energy-preserving}
\end{equation}
Equation \eqref{eq:energy-preserving} is a strong and deliberate
constraint. It says the interaction can only connect
\emph{degenerate} eigenstates of $H_0$ --- it shuffles excitations among
qubits without changing the total bare energy. We will see in
Sec.~\ref{sec:interaction} that this single requirement, together with a
resonance condition on the gaps, \emph{uniquely} dictates the form of
the collector's three--body term.

\begin{approxbox}[(A5) energy--preserving interaction, \eqref{eq:energy-preserving}]
\textbf{Where:} $[H_{\mathrm{int}},H_0]=0$. \textbf{Why:} in the weak
internal--coupling regime, any non--commuting (energy--changing) part of
the interaction is off--resonant and averages away under the same
rotating--wave logic as (A3); only the energy--conserving, secular part
survives to influence the steady state. \textbf{When:} coupling
$\chi\ll$ the gaps, so that off--resonant transitions are suppressed by
$\chi/\epsilon$. This is what justifies the local master equation
\eqref{eq:local-me} below.
\end{approxbox}

\subsection{The local master equation}
Each qubit $k$ is weakly coupled to its own bath $B_k$ at inverse
temperature $\beta_k$. Under weak coupling the machine obeys the
\emph{local} master equation
\begin{equation}
\dot\rho = -\ii[H_0 + H_{\mathrm{int}},\rho]
+ \sum_k \calL^{(k)}[\rho],
\label{eq:local-me}
\end{equation}
with each $\calL^{(k)}$ the reset dissipator \eqref{eq:reset} for its own
bath. This is [paper Eq.~(3)]. ``Local'' means each dissipator is built
from the \emph{bare} qubit operators, not the dressed eigenstates of
$H_0+H_{\mathrm{int}}$.

\begin{approxbox}[(A6) local vs.\ global master equation]
\textbf{Where:} the dissipators in \eqref{eq:local-me} act on individual
bare qubits. \textbf{Why:} when $H_{\mathrm{int}}$ is weak compared with
the gaps, the eigenstates of $H_0+H_{\mathrm{int}}$ are close to the bare
product states, and the local dissipator is both simpler and
thermodynamically consistent. \textbf{When:} $\chi \ll \epsilon_k$; in
the opposite (strong internal coupling) regime one must use the global
dissipator built from the dressed eigenoperators, otherwise small
violations of the second law can appear (Hofer \emph{et al.}, cited as
[51] in the paper). We stay local throughout.
\end{approxbox}

\subsection{Local detailed balance and heat currents}
The key property inherited from Sec.~\ref{sec:open} is that each local
dissipator has the local Gibbs state as a zero,
\begin{equation}
\calL^{(k)}[\tau(\beta_k)] = 0 ,
\label{eq:ldb}
\end{equation}
which is [paper Eq.~(4)]. The heat current into bath $k$ is again
\eqref{eq:current-def}, and, crucially, in a nonequilibrium steady state
(NESS) the machine's own entropy is constant, so all the entropy the
machine produces is dumped into the baths.

\subsection{Entropy production and the second law}
Define the machine's von Neumann entropy $S(\rho) = -\Tr[\rho\ln\rho]$.
The entropy production rate is
\begin{equation}
\dot\Sigma := \dot S(\rho(t)) - \sum_k \beta_k\, j_k(t) \;\ge\; 0,
\label{eq:entropy-prod}
\end{equation}
which is [paper Eq.~(7)] and is the second law for this class of
machines. In the steady state $\dot S = 0$, so
$\dot\Sigma = -\sum_k\beta_k j_k \ge 0$ is purely the weighted heat
dumped into the environments. This quantity is the ``dissipation'' whose
trade--off against computational accuracy we quantify in
Part~\ref{part:notgate}. It is also the microscopic meaning of the
``information lost to the environment'' discussed qualitatively in the
Floquet--buffer note.

\begin{keybox}
\textbf{End of Part I.} We now have every generic tool: a free
Hamiltonian $H_0=\sum_k\epsilon_k n_k$, an energy--preserving interaction
$H_{\mathrm{int}}$ (still to be constructed), a local reset dissipator
$\sum_k\calL^{(k)}$ with detailed balance, a heat current
$j_k=\gamma_k\epsilon_k[g(\beta_k\epsilon_k)-p_k]$, and an entropy
production $\dot\Sigma\ge0$. Part II uses two qubits to build the virtual
qubit; Part III assembles the NOT gate.
\end{keybox}

% #####################################################################
\part{The virtual qubit and the resonant interaction}
\label{part:virtual}
% #####################################################################

\section{Two thermal qubits and the virtual qubit}
\label{sec:virtual}

\subsection{The product thermal state}
Take two qubits $C_0,C_1$ with gaps $\epsilon_0$ and
$\epsilon_1\le\epsilon_0$, each in contact with its own bath at inverse
temperature $\beta_0,\beta_1$. In the absence of any interaction, each
qubit independently reaches its Gibbs state \eqref{eq:gibbs-explicit},
\begin{equation}
\tau_{C_i}(\beta_i)
= \frac{1}{Z_{C_i}}\Big(\ket{0}\bra{0}_{C_i}
+ \ee^{-\beta_i\epsilon_i}\ket{1}\bra{1}_{C_i}\Big),\quad
Z_{C_i}=1+\ee^{-\beta_i\epsilon_i},
\label{eq:tauCi}
\end{equation}
[paper Eq.~(8)]. The joint state is the product
$\tau_{C_0}(\beta_0)\otimes\tau_{C_1}(\beta_1)$, diagonal in the four
states $\ket{i}_{C_0}\ket{j}_{C_1}$, $i,j\in\{0,1\}$, with populations
\begin{equation}
\begin{aligned}
&P_{00}=p_0(0)\,p_1(0), &\quad& P_{01}=p_0(0)\,p_1(1),\\
&P_{10}=p_0(1)\,p_1(0), &\quad& P_{11}=p_0(1)\,p_1(1),
\end{aligned}
\label{eq:Pij}
\end{equation}
where $P_{ij}:=\bra{ij}\,\tau_{C_0}\!\otimes\tau_{C_1}\,\ket{ij}$ is the
joint population of $\ket{i}_{C_0}\ket{j}_{C_1}$, and
$p_i(1)=g(\beta_i\epsilon_i)$, $p_i(0)=1-p_i(1)$.

\subsection{The virtual qubit subspace}
Focus on the two ``crossed'' states
\begin{equation}
\ket{0}_v := \ket{0}_{C_0}\ket{1}_{C_1},\qquad
\ket{1}_v := \ket{1}_{C_0}\ket{0}_{C_1}.
\label{eq:virtual-states}
\end{equation}
This is [paper Eq.~(15)] (collector convention). They span a
two--dimensional subspace of the four--dimensional joint space --- a
\emph{virtual qubit}. Its energy gap is the difference of the bare
energies of the two states,
\begin{equation}
\epsilon_v := E(\ket{1}_v) - E(\ket{0}_v) = \epsilon_0 - \epsilon_1 .
\label{eq:eps-v}
\end{equation}
The name ``virtual'' is literal: no physical two--level system with gap
$\epsilon_v$ exists; it is an effective qubit encoded in a corner of the
joint Hilbert space of $C_0$ and $C_1$.

\subsection{The virtual temperature}
Assign the virtual qubit a temperature exactly as we did for a real one,
via its population ratio \eqref{eq:db-ratio}. Using
\eqref{eq:Pij} and the Boltzmann ratios
$p_i(1)/p_i(0)=\ee^{-\beta_i\epsilon_i}$,
\begin{equation}
\ee^{-\beta_v\epsilon_v}
:= \frac{\bra{1}_v\,\tau_{C_0}\!\otimes\tau_{C_1}\,\ket{1}_v}
        {\bra{0}_v\,\tau_{C_0}\!\otimes\tau_{C_1}\,\ket{0}_v}
= \frac{P_{10}}{P_{01}}
= \frac{p_0(1)p_1(0)}{p_0(0)p_1(1)}
= \frac{\ee^{-\beta_0\epsilon_0}}{\ee^{-\beta_1\epsilon_1}} .
\label{eq:betav-ratio}
\end{equation}
This is [paper Eq.~(10)]. Taking logarithms,
$\beta_v\epsilon_v = \beta_0\epsilon_0-\beta_1\epsilon_1$, and dividing
by $\epsilon_v=\epsilon_0-\epsilon_1$,
\begin{equation}
\boxed{\;
\beta_v = \frac{\epsilon_0}{\epsilon_0-\epsilon_1}\,\beta_0
        - \frac{\epsilon_1}{\epsilon_0-\epsilon_1}\,\beta_1 .
\;}
\label{eq:betav}
\end{equation}
This is [paper Eq.~(11)/(16)]. The virtual (inverse) temperature is a
\emph{linear combination of the input inverse temperatures}, with
weights set by the energy gaps. \emph{This is the weighted sum of the
perceptron}, and it is the reason the machine computes. Note that
$\beta_v$ can exceed both $\beta_0,\beta_1$, lie between them, or go
negative --- the three regimes below.

\subsection{Refrigerator, heat pump, heat engine}
Holding $\beta_0$ fixed and varying $\beta_1$ in \eqref{eq:betav}:
\begin{itemize}[leftmargin=*]
\item $\beta_1<\beta_0$ (hot input): $\beta_v$ is larger than both
$\beta_0$ and $\beta_1$ --- the virtual qubit is \emph{colder} than
either bath; coupling a target to it \emph{cools} the target
(refrigerator).
\item $\beta_0<\beta_1<(\epsilon_0/\epsilon_1)\beta_0$: $\beta_v$ is
smaller than both --- \emph{hotter} virtual qubit (heat pump).
\item $\beta_1>(\epsilon_0/\epsilon_1)\beta_0$: $\beta_v<0$ ---
population inversion (heat engine).
\end{itemize}
These are the regions of [paper Fig.~2]. The ability of a single knob
$\beta_1$ to move the target across cooling/heating is the physical basis
of temperature inversion, i.e.\ of the NOT gate.

\begin{implbox}[virtual temperature]
Directly: \texttt{beta\_v = (eps0/(eps0-eps1))*beta0 -
(eps1/(eps0-eps1))*beta1}. As a check, build
\texttt{tau01 = tensor(thermal(beta0,eps0), thermal(beta1,eps1))} and
verify \eqref{eq:betav-ratio}: the ratio of the $\ket{10}$ and
$\ket{01}$ diagonal entries equals $\ee^{-\beta_v(\epsilon_0-\epsilon_1)}$.
\end{implbox}

% ---------------------------------------------------------------------
\section{Adding the target qubit: the resonant three--body interaction}
\label{sec:interaction}

We now derive --- not quote --- the collector's interaction Hamiltonian,
the object your reader most wants to see written out and justified.

\subsection{The resonance condition}
Introduce a third qubit $C_z$ with gap $\epsilon_z$, intended to be
placed in ``thermal contact'' with the virtual qubit so that it inherits
$\beta_v$. For a weak, energy--preserving interaction to move excitations
between the virtual qubit and $C_z$, the two must be \emph{resonant}:
the energy released by de--exciting the virtual qubit must exactly
excite $C_z$. Hence we \emph{choose}
\begin{equation}
\boxed{\;\epsilon_z = \epsilon_v = \epsilon_0 - \epsilon_1 .\;}
\label{eq:resonance}
\end{equation}
This is a design equation, imposed by hand on the fabricated gaps
[paper, above Eq.~(14)]. Its consequence: the two joint states
\begin{equation}
\ket{1}_v\ket{0}_{C_z}=\ket{1}_{C_0}\ket{0}_{C_1}\ket{0}_{C_z}
=\ket{100},\qquad
\ket{0}_v\ket{1}_{C_z}=\ket{0}_{C_0}\ket{1}_{C_1}\ket{1}_{C_z}
=\ket{011},
\end{equation}
become \emph{degenerate}. Their bare energies are
\begin{equation}
E(\ket{100}) = \epsilon_0,\qquad
E(\ket{011}) = \epsilon_1 + \epsilon_z
\overset{\eqref{eq:resonance}}{=}\epsilon_0 .
\end{equation}
(Ordering of kets is $(C_0C_1C_z)$ throughout.)

\subsection{Deriving the interaction from \texorpdfstring{$[H_{\mathrm{int}},H_0]=0$}{[Hint,H0]=0}}
The energy--preserving requirement \eqref{eq:energy-preserving} says
$H_{\mathrm{int}}$ may only connect degenerate eigenstates of $H_0$.
Among all three--qubit basis states, the single resonant pair created by
\eqref{eq:resonance} is $\{\ket{100},\ket{011}\}$. The most general
Hermitian operator connecting \emph{only} this pair is
\begin{equation}
\boxed{\;
H_{\mathrm{int}} = \chi\big(\ket{100}\bra{011}
+ \ket{011}\bra{100}\big),
\;}
\label{eq:Hint-kets}
\end{equation}
with a single real coupling $\chi$ (a phase on $\chi$ can be absorbed
into the definition of the states). Written with Pauli operators, using
$\ket{100}\bra{011}=\sigma_+^{(0)}\sigma_-^{(1)}\sigma_-^{(z)}$,
\begin{equation}
H_{\mathrm{int}} = \chi\Big(
\sigma_+^{(0)}\sigma_-^{(1)}\sigma_-^{(z)}
+ \sigma_-^{(0)}\sigma_+^{(1)}\sigma_+^{(z)}\Big),
\label{eq:Hint-pauli}
\end{equation}
or, in the mixed notation of [paper Eq.~(14)],
\begin{equation}
H_{\mathrm{int}} = \chi\,\ket{1}\bra{0}_{C_0}\otimes
\ket{0}\bra{1}_{C_1}\otimes\ket{0}\bra{1}_{C_z} + \mathrm{h.c.}
\label{eq:Hint-paper}
\end{equation}
Equations \eqref{eq:Hint-kets}--\eqref{eq:Hint-paper} are one and the
same operator; \eqref{eq:Hint-paper} is verbatim the paper's. One checks
directly that $[H_{\mathrm{int}},H_0]=0$ because it connects only
degenerate states.

\subsection{Why this term, and why not the others}
Five questions, five answers.
\begin{enumerate}[label=(\roman*),leftmargin=*]
\item \textbf{Why does energy conservation force exactly this pair?}
Because \eqref{eq:resonance} makes $\ket{100}$ and $\ket{011}$ the
\emph{only} degenerate pair differing by a single allowed excitation
rearrangement; any other off--diagonal term would connect states of
different bare energy and violate \eqref{eq:energy-preserving}.
\item \textbf{Why three--body and not pairwise?} The virtual qubit
\eqref{eq:virtual-states} is a \emph{joint} property of $C_0$ and $C_1$.
A two--body term $\sigma_+^{(0)}\sigma_-^{(1)}$ merely rotates the
excitation between $C_0$ and $C_1$ (an internal rotation of the virtual
qubit) and transfers nothing to $C_z$. To let $C_z$ ``read'' the virtual
qubit you must flip all three spins at once --- an irreducibly
three--body operator.
\item \textbf{What does it do physically?} Rewrite
\eqref{eq:Hint-kets} in virtual--qubit variables:
\begin{equation}
H_{\mathrm{int}} = \chi\big(\ket{1}_v\ket{0}_{C_z}
\bra{0}_v\bra{1}_{C_z} + \mathrm{h.c.}\big),
\label{eq:Hint-swap}
\end{equation}
[paper Eq.~(12)] --- a resonant \textsc{swap} (beam splitter) between the
virtual qubit and $C_z$. \emph{This is the microscopic reason $C_z$
thermalizes to $\beta_v$}: the interaction literally swaps the virtual
qubit's population into $C_z$, so in steady state $C_z$ carries the
virtual population, i.e.\ the virtual temperature.
\item \textbf{Why drop the counter--rotating / non--resonant terms?}
Operators like $\sigma_+^{(0)}\sigma_+^{(1)}\sigma_+^{(z)}$ create three
excitations and oscillate at $\epsilon_0+\epsilon_1+\epsilon_z$ in the
interaction picture; being far off resonance they average to zero on the
relaxation time (the same RWA as (A3)/(A5)), with error $O(\chi/\epsilon)$.
\item \textbf{Why is there no bare $C_0$--$C_1$ coupling or
$\sigma_z\sigma_z$ term?} Such terms would hybridize the two input
qubits and spoil the \emph{product} thermal structure
$\tau_{C_0}\otimes\tau_{C_1}$ from which $\beta_v$ was \emph{defined} in
\eqref{eq:betav-ratio}. They are deliberately absent. The three lines a
schematic draws between $C_0,C_1,C_z$ represent the single three--body
term \eqref{eq:Hint-kets}, not three independent couplings.
\end{enumerate}

\begin{remark}[a frequent coding error]
The resonant pair is $\ket{100}\leftrightarrow\ket{011}$. A tempting but
\emph{wrong} choice is $\ket{101}\leftrightarrow\ket{010}$, whose bare
energies are $\epsilon_0+\epsilon_z$ and $\epsilon_1$; these differ by
$\epsilon_0+\epsilon_z-\epsilon_1 = 2\epsilon_z\neq0$. That term is
off--resonant and, at $\chi\ll\epsilon_z$, dynamically inert --- the
collector would not work. Always build $H_{\mathrm{int}}$ from the
degenerate pair fixed by \eqref{eq:resonance}.
\end{remark}

\begin{implbox}[the collector Hamiltonian]
With QuTiP and ordering $(C_0C_1C_z)$:
\begin{quote}\ttfamily
n0,n1,nz = num on each site (tensor with identities)\\
H0 = eps0*n0 + eps1*n1 + eps\_z*nz\\
k100 = tensor(basis(2,1),basis(2,0),basis(2,0))\\
k011 = tensor(basis(2,0),basis(2,1),basis(2,1))\\
Hint = chi*(k100*k011.dag() + k011*k100.dag())\\
H = H0 + Hint
\end{quote}
Verify $\|[H_0,H_{\mathrm{int}}]\|=0$ before trusting anything: it is the
one--line check that the resonance \eqref{eq:resonance} was implemented
correctly.
\end{implbox}

\subsection{Steady state of the swap: \texorpdfstring{$C_z\to\beta_v$}{Cz to betav}}
With $C_z$ additionally reset--coupled (rate $\mu$) to a bath $B_z$ at
$\beta_z$, the competition is between the swap
\eqref{eq:Hint-swap}, which pulls $C_z$ toward $\beta_v$, and the bath,
which pulls it toward $\beta_z$. Under the time--scale separation
introduced next (Part III), the current the collector delivers to $B_z$
is [paper Eq.~(17)]
\begin{equation}
\boxed{\; j_{\calC} = \mu\,\epsilon_z\big[\gz(\beta_z) - \gz(\beta_v)\big],
\qquad \gz(x):=g(x\,\epsilon_z). \;}
\label{eq:jC}
\end{equation}
When $\beta_v>\beta_z$ the collector cools $B_z$ (refrigerator); when
$\beta_v<\beta_z$ it heats it (heat pump). Setting $j_{\calC}=0$ in
isolation gives $\beta_z=\beta_v$: \emph{the collector alone drives the
output reservoir to the virtual temperature}. That linearity
(cf.\ \eqref{eq:betav}) is precisely why the collector by itself is only
an inverting \emph{linear amplifier} --- the motivation for the
modulator, taken up in Part III.

\begin{keybox}
\textbf{End of Part II.} The collector's Hamiltonian is
$H_{\calC}=H_0+H_{\mathrm{int}}$ with
$H_0=\epsilon_0 n_0+\epsilon_1 n_1+\epsilon_z n_z$ and the resonant
three--body swap \eqref{eq:Hint-kets}. It thermalizes $C_z$ (and hence
$B_z$) to the virtual temperature $\beta_v$ of \eqref{eq:betav}, a
weighted sum of the inputs --- the perceptron pre--activation. Part III
adds the modulator to make the response nonlinear and reads out the bit.
\end{keybox}

% #####################################################################
\part{The NOT gate: collector, modulator, dynamics, logic}
\label{part:notgate}
% #####################################################################

\section{The collector alone is a linear amplifier}
\label{sec:collector-linear}
Assemble the collector: three qubits $C_0,C_1,C_z$ with the free
Hamiltonian $H_0=\epsilon_0 n_0+\epsilon_1 n_1+\epsilon_z n_z$, the
resonant three--body interaction \eqref{eq:Hint-kets}, and three local
reset dissipators (rates $\gamma$ for $C_0,C_1$; $\mu$ for $C_z$). Its
master equation is [paper Eq.~(13)]
\begin{equation}
\dot\rho_{\calC} = -\ii[H_0+H_{\mathrm{int}},\rho_{\calC}]
+ \calL^{(0)}[\rho_{\calC}] + \calL^{(1)}[\rho_{\calC}]
+ \calL^{(z)}[\rho_{\calC}].
\label{eq:collector-me}
\end{equation}
As shown in Sec.~\ref{sec:interaction}, the swap thermalizes $C_z$ toward
the virtual temperature $\beta_v$, delivering to $B_z$ the current
\eqref{eq:jC}, $j_{\calC}=\mu\epsilon_z[\gz(\beta_z)-\gz(\beta_v)]$. In
isolation the fixed point $j_{\calC}=0$ is $\beta_z=\beta_v$, and by
\eqref{eq:betav} this is an \emph{affine} (linear) function of the input:
\begin{equation}
\beta_z^{\text{collector only}} = \beta_v
= \frac{\epsilon_0}{\epsilon_z}\,\beta_0
- \frac{\epsilon_1}{\epsilon_z}\,\beta_1 .
\label{eq:linear-response}
\end{equation}
The output inverse temperature is a straight line in $\beta_1$ with
negative slope: an \emph{inverting linear amplifier}.

\section{Why we need the modulator}
\label{sec:why-modulator}
Here is the answer to the question ``why do we even need the
modulator?''. A linear amplifier is a \emph{bad} logic gate for one
concrete reason: it does not restore levels. Suppose the input carries a
small fluctuation $\beta_1\to\beta_1+\delta$. By
\eqref{eq:linear-response} the output moves by
$-(\epsilon_1/\epsilon_z)\,\delta$; with $\epsilon_1/\epsilon_z$ made
large (needed for a steep, decisive gate) the \emph{noise is amplified}
along with the signal. Cascading such gates amplifies noise without
bound, causing bit flips. A good gate must instead \emph{saturate}: for
inputs safely inside a logical region the output should barely move
(flat response, noise rejected), while it should swing sharply when the
input crosses the decision boundary. That is precisely a
\emph{sigmoidal}, not linear, transfer curve.

\begin{keybox}
\textbf{The role of the modulator (one sentence).} The collector
supplies the perceptron's \emph{weighted sum} $\beta_v$
\eqref{eq:betav}; the modulator supplies the perceptron's \emph{nonlinear
activation} by confining the output to a fixed logical window
$[\beta_{\mathrm{hot}},\beta_{\mathrm{cold}}]$ and thereby saturating the
response. Without it the ``activation function'' is the identity
$f(y)=y$, which is known to make a useless classifier.
\end{keybox}

\section{The modulator: one qubit, two baths}
\label{sec:modulator}
The modulator $\calM$ is a \emph{single} qubit $C_\calM$ with gap chosen
resonant with the output channel,
\begin{equation}
H_\calM = \epsilon_\calM\, n_\calM ,\qquad \epsilon_\calM = \epsilon_z ,
\end{equation}
placed in simultaneous contact with \emph{two} baths: a reference bath
$B_r$ at fixed $\beta_r$ (rate $\gamma$) and the output bath $B_z$ (rate
$\mu'$). Its state obeys [paper Eq.~(18)]
\begin{equation}
\dot\rho_\calM = \calL^{(r)}[\rho_\calM] + \calL^{(z)}[\rho_\calM]
= \gamma\big(\tau(\beta_r)-\rho_\calM\big)
+ \mu'\big(\tau(\beta_z)-\rho_\calM\big).
\label{eq:mod-me}
\end{equation}
Both dissipators are diagonal, so it suffices to track the excited
population $p_\calM:=\bra{1}\rho_\calM\ket{1}$:
\begin{equation}
\dot p_\calM = \gamma\big[\gz(\beta_r)-p_\calM\big]
+ \mu'\big[\gz(\beta_z)-p_\calM\big].
\label{eq:mod-pop}
\end{equation}
Setting $\dot p_\calM=0$,
\begin{equation}
p_\calM^{\mathrm{ss}}
= \frac{\gamma\,\gz(\beta_r)+\mu'\,\gz(\beta_z)}{\gamma+\mu'}
\;\xrightarrow{\ \mu'\ll\gamma\ }\; \gz(\beta_r).
\label{eq:mod-ss}
\end{equation}

\begin{approxbox}[(A7) modulator pinned to the reference, \eqref{eq:mod-ss}]
\textbf{Where:} we set $\mu'\ll\gamma$, so $p_\calM^{\mathrm{ss}}\approx
\gz(\beta_r)$. \textbf{Why:} the strong reference coupling holds the
modulator qubit at $\beta_r$, essentially unperturbed by the (weak)
output coupling. \textbf{When:} $\mu'/\gamma\to0$; corrections are
$O(\mu'/\gamma)$. This is what makes the modulator a stiff ``anchor''
whose only slow action is to bleed a controlled current into $B_z$.
\end{approxbox}

The steady--state heat current the modulator delivers to $B_z$ is
\eqref{eq:current-reset} with $p_\calM=\gz(\beta_r)$, [paper Eq.~(19)]
\begin{equation}
\boxed{\; j_\calM = \mu'\,\epsilon_z\big[\gz(\beta_z)-\gz(\beta_r)\big]. \;}
\label{eq:jM}
\end{equation}
The modulator therefore continuously pulls $\beta_z$ toward $\beta_r$
with strength $\mu'$: a single--qubit \emph{heat valve} that competes
with the collector current \eqref{eq:jC}.

\section{Dynamics of the full machine: two time scales}
\label{sec:dynamics}
The output reservoir $B_z$ has finite heat capacity $\calC$, so its
temperature $T_z=1/\beta_z$ moves under the total incoming current
$j_{\calC}+j_\calM$ (calorimetric law $\calC\dot T_z=j_{\calC}+j_\calM$).
In terms of $\beta_z=1/T_z$ (units $\kB=1$),
\begin{equation}
\boxed{\;
\dot\beta_z = -\frac{\beta_z^2}{\calC}\big(j_{\calC}+j_\calM\big).
\;}
\label{eq:calorimetric}
\end{equation}
This is [paper Eq.~(20)]. The consistency of inserting the
\emph{steady--state} currents \eqref{eq:jC}, \eqref{eq:jM} into
\eqref{eq:calorimetric} rests on a separation of time scales built into
the design,
\begin{equation}
\gamma \;\gg\; \mu,\mu' .
\label{eq:timescales}
\end{equation}

\begin{approxbox}[(A8) adiabatic elimination / time--scale separation, \eqref{eq:timescales}]
\textbf{Where:} we use the \emph{steady--state} $C_z\to\beta_v$ and
$C_\calM\to\beta_r$ inside the slow equation \eqref{eq:calorimetric}.
\textbf{Why:} with $\gamma\gg\mu,\mu'$ the internal qubits relax fast
compared with the drift of $\beta_z(t)$, so at each instant they are
slaved to the current value of $\beta_z$. \textbf{When:} $\mu,\mu'\ll
\gamma$; the error is $O(\mu/\gamma)$. This is the autonomous analogue of
the quasi--static assumption in circuit analysis, and it is what lets us
treat $\beta_z(t)$ as the single slow logical variable.
\end{approxbox}

\section{Steady state and the transfer characteristic}
\label{sec:transfer}
The output settles when $\dot\beta_z=0$, i.e.\ $j_{\calC}+j_\calM=0$.
Substituting \eqref{eq:jC} and \eqref{eq:jM},
\begin{equation}
\mu\big[\gz(\beta_z)-\gz(\beta_v)\big]
+ \mu'\big[\gz(\beta_z)-\gz(\beta_r)\big]=0,
\end{equation}
so, with the branching ratio $\Delta:=\mu/(\mu+\mu')$,
\begin{equation}
\boxed{\;
\gz(\beta_z^\infty) = \Delta\,\gz(\beta_v)+(1-\Delta)\,\gz(\beta_r).
\;}
\label{eq:ss-condition}
\end{equation}
This is [paper Eq.~(21)]. The output excitation is a convex combination
of the collector's target ($\beta_v$) and the modulator's anchor
($\beta_r$), weighted purely by coupling rates.

\subsection{Fixing the modulator: confinement to the logical window}
For $\beta_z^\infty$ to be a valid logical signal we demand it stay in a
window $[\beta_{\mathrm{hot}},\beta_{\mathrm{cold}}]$
($\beta_{\mathrm{cold}}>\beta_{\mathrm{hot}}$) even when the virtual
temperature saturates. Since $\gz$ is monotonically decreasing,
$\gz(\beta_v)\to0$ as $\beta_v\to+\infty$ and $\gz(\beta_v)\to1$ as
$\beta_v\to-\infty$. Impose the two rail conditions on
\eqref{eq:ss-condition}:
\begin{align}
\beta_v\to+\infty:&\quad (1-\Delta)\gz(\beta_r)
\overset{!}{=}\gz(\beta_{\mathrm{cold}}),
\label{eq:rail1}\\
\beta_v\to-\infty:&\quad \Delta+(1-\Delta)\gz(\beta_r)
\overset{!}{=}\gz(\beta_{\mathrm{hot}}).
\label{eq:rail2}
\end{align}
Subtracting \eqref{eq:rail1} from \eqref{eq:rail2} gives the two design
equations that fix the modulator's free parameters:
\begin{equation}
\boxed{\;
\Delta = \gz(\beta_{\mathrm{hot}})-\gz(\beta_{\mathrm{cold}}),
\qquad
\gz(\beta_r)=\frac{\gz(\beta_{\mathrm{cold}})}{1-\Delta}. \;}
\label{eq:mod-design}
\end{equation}
These pin $(\mu'/\mu,\beta_r)$ from the chosen logical window. This is the
main--text version of [paper Supplementary Material A]; the sign
conventions there ($\beta_{\min},\beta_{\max},\kappa$) are identical
after relabelling.

\subsection{The transfer function}
Substituting \eqref{eq:mod-design} back into \eqref{eq:ss-condition},
$(1-\Delta)\gz(\beta_r)=\gz(\beta_{\mathrm{cold}})$ and
$\Delta=\gz(\beta_{\mathrm{hot}})-\gz(\beta_{\mathrm{cold}})$, so
\begin{align}
\gz(\beta_z^\infty)
&= \gz(\beta_{\mathrm{hot}})\gz(\beta_v)
+ \gz(\beta_{\mathrm{cold}})\big(1-\gz(\beta_v)\big) =: Q(\beta_v).
\label{eq:Qdef}
\end{align}
Inverting $\gz(\beta)=(1+\ee^{\beta\epsilon_z})^{-1}$
(so $\beta=\epsilon_z^{-1}\ln[\gz^{-1}-1]$) yields the exact transfer
characteristic [paper Eq.~(22)]
\begin{equation}
\boxed{\;
\beta_z^\infty = \frac{1}{\epsilon_z}\,
\ln\!\big[Q(\beta_v)^{-1}-1\big],
\quad
Q(\beta_v)=\gz(\beta_{\mathrm{hot}})\gz(\beta_v)
+\gz(\beta_{\mathrm{cold}})\big(1-\gz(\beta_v)\big). \;}
\label{eq:transfer}
\end{equation}

\begin{implbox}[transfer curve]
Compute \texttt{beta\_v} from \eqref{eq:betav}, then \texttt{gz(x)=1/(1+exp(x*eps\_z))},
\texttt{Q = gz(bhot)*gz(bv)+gz(bcold)*(1-gz(bv))}, and
\texttt{bz\_inf = log(1/Q - 1)/eps\_z}. Sweeping $\beta_1$ reproduces
[paper Fig.~3B]. As a cross--check, integrate \eqref{eq:calorimetric}
with a stiff ODE solver and confirm the fixed point matches
\eqref{eq:transfer}.
\end{implbox}

\section{Emergence of the sigmoid}
\label{sec:sigmoid}
The link to the perceptron activation is now exact. First note the
identity, from $\beta_v\epsilon_z=\beta_0\epsilon_0-\beta_1\epsilon_1$
(multiply \eqref{eq:betav} by $\epsilon_z=\epsilon_0-\epsilon_1$),
\begin{equation}
\gz(\beta_v)=\frac{1}{1+\ee^{\beta_v\epsilon_z}}
= f\big(\beta_0\epsilon_0-\beta_1\epsilon_1\big),\qquad
f(x):=\frac{1}{1+\ee^{x}},
\label{eq:exact-sigmoid}
\end{equation}
with \emph{no approximation}: the collector's virtual population is
already a sigmoid of the perceptron pre--activation
$x=\beta_0\epsilon_0-\beta_1\epsilon_1$. The modulator's job is only to
map it affinely onto the temperature axis. Expanding \eqref{eq:transfer}
for small $\epsilon_z$ using
$\gz(\beta)=\tfrac12-\tfrac{\beta\epsilon_z}{4}+O(\epsilon_z^3)$:
\begin{equation}
Q(\beta_v)=\tfrac12-\frac{\epsilon_z}{4}\,B+O(\epsilon_z^2),\quad
B:=\beta_{\mathrm{hot}}\gz(\beta_v)
+\beta_{\mathrm{cold}}\big(1-\gz(\beta_v)\big),
\end{equation}
so that $\frac{1-Q}{Q}=1+\epsilon_z B+O(\epsilon_z^2)$ and
\begin{equation}
\beta_z^\infty=\frac1{\epsilon_z}\ln\frac{1-Q}{Q}
= B+O(\epsilon_z)
= \beta_{\mathrm{cold}}
+(\beta_{\mathrm{hot}}-\beta_{\mathrm{cold}})\,\gz(\beta_v)+O(\epsilon_z).
\label{eq:activation}
\end{equation}
Combining \eqref{eq:exact-sigmoid} and \eqref{eq:activation},
\begin{equation}
\boxed{\;
\beta_z^\infty \simeq \beta_{\mathrm{cold}}
+(\beta_{\mathrm{hot}}-\beta_{\mathrm{cold}})\,
f\big(\beta_0\epsilon_0-\beta_1\epsilon_1\big),
\;}
\label{eq:sigmoid-final}
\end{equation}
which is [paper Eq.~(23)]: a sigmoidal interpolation between the two
logical rails whose argument is exactly the perceptron pre--activation.
As $\epsilon_1\to\infty$ the sigmoid sharpens into the ideal inverted
step (perfect NOT), at the cost of dissipation (Sec.~\ref{sec:tradeoff}).

\begin{approxbox}[(A9) small--$\epsilon_z$ (sigmoid) expansion, \eqref{eq:activation}]
\textbf{Where:} we expand $\gz$ and the logarithm to first order in
$\epsilon_z$. \textbf{Why:} for $\epsilon_z\ll\epsilon_1$ the output gap
is the smallest scale and the transfer curve is dominated by its linear
term, which is a clean sigmoid. \textbf{When:} $\epsilon_z\ll
1/\beta_{\mathrm{hot,cold}},\epsilon_1$. For larger $\epsilon_z$ the
curve still inverts but the parameter roles ($\beta_0$ as threshold,
$\epsilon_0$ as steepness) blur.
\end{approxbox}

\subsection{Threshold and gain}
Differentiating \eqref{eq:sigmoid-final} with $f'(x)=-f(x)(1-f(x))$ and
$x=\beta_0\epsilon_0-\beta_1\epsilon_1$ (so $\dd x/\dd\beta_1=-\epsilon_1$),
\begin{equation}
\mathcal{A}(\beta_1):=\dv{\beta_z^\infty}{\beta_1}
= -\,\Delta\beta\,\epsilon_1\,f(x)\big(1-f(x)\big),\qquad
\Delta\beta:=\beta_{\mathrm{cold}}-\beta_{\mathrm{hot}}.
\end{equation}
The switching threshold is where $x=0$, i.e.\
$\beta_1^\ast=\beta_0\epsilon_0/\epsilon_1
=\beta_0(1+\epsilon_z/\epsilon_1)\approx\beta_0$; there $f(1-f)=1/4$ and
\begin{equation}
\mathcal{A}_{\max}=-\frac{\epsilon_1\,\Delta\beta}{4},
\label{eq:gain}
\end{equation}
recovering [paper Supp.~A, Eq.~(A10)] (for $\beta_{\mathrm{hot}}=0,
\beta_{\mathrm{cold}}=1$, $\mathcal{A}=-\epsilon_0/4+O(\epsilon_z^2)$
with $\epsilon_0\approx\epsilon_1$). \textbf{This makes the trade--off
visible already:} steeper gate $\Rightarrow$ larger $\epsilon_1$
$\Rightarrow$ (Sec.~\ref{sec:tradeoff}) more dissipation. The qubit gap
$\epsilon_1$ is the ``transconductance'' of the device.

\section{Logic encoding and readout error}
\label{sec:logic}
\subsection{Encoding and decoding}
Logical values are placed on, and read from, temperatures via [paper
Eq.~(24)]
\begin{equation}
x=\begin{cases}0,&\beta_1=\beta_{\mathrm{hot}}\\
1,&\beta_1=\beta_{\mathrm{cold}}\end{cases}
\qquad
y=\begin{cases}
0,&\beta_z^\infty\le(1+\delta)\beta_{\mathrm{hot}}\\
1,&\beta_z^\infty\ge(1-\delta)\beta_{\mathrm{cold}}\\
\varnothing,&\text{otherwise,}\end{cases}
\label{eq:encoding}
\end{equation}
where $\delta$ is a tolerance (guard band) and $\varnothing$ is an
invalid readout. Inversion follows from the sign structure of $\beta_v$:
a hot input ($\beta_1=\beta_{\mathrm{hot}}<\beta_0$) makes the collector
a refrigerator, $\gz(\beta_v)\to0$, so
$\beta_z^\infty\to\beta_{\mathrm{cold}}$ ($y=1$); a cold input makes it a
heat pump, $\gz(\beta_v)\to1$, $\beta_z^\infty\to\beta_{\mathrm{hot}}$
($y=0$). That is the NOT gate.

\subsection{The Gaussian readout--error model}
Because the machine is stochastic, the finite bath fluctuates around
$\beta_z^\infty$. Model the actual response as Gaussian [paper
Eqs.~(25)--(26)],
\begin{equation}
T(\beta_z\mid\beta_1)\propto \mathcal N\big(\beta_z^\infty,\;C\big),
\end{equation}
justified by the central limit theorem for the many--particle reservoir
(larger heat capacity $C$ $\Rightarrow$ wider fluctuations). With input
distribution $p(x)$, encoding $e(\beta_1\mid x)$ and decoding
$d(y\mid\beta_z)$ from \eqref{eq:encoding}, the channel is
\begin{equation}
p(y\mid x)=\int d(y\mid\beta_z)\,T(\beta_z\mid\beta_1)\,
e(\beta_1\mid x)\,\dd\beta_1\,\dd\beta_z,
\end{equation}
and the average error is [paper Eq.~(27)]
\begin{equation}
\langle\xi\rangle=\sum_{x\in\{0,1\}}\sum_{y\in\{0,1\}}
p(x)\,p(y\mid x)\,\delta_{xy}
\quad(\text{probability output}\neq\text{desired}).
\label{eq:error}
\end{equation}
The sharper the transfer curve (larger $\epsilon_1$), the smaller
$\langle\xi\rangle$ --- the gate quality improves with energy invested.

\section{The error--dissipation trade--off}
\label{sec:tradeoff}
The dissipation is the entropy production \eqref{eq:entropy-prod}
integrated over the run. In the slow regime the machine's own entropy is
constant, so the total rate splits into collector and modulator pieces,
\begin{equation}
\dot\Sigma=\dot\Sigma_{\calC}+\dot\Sigma_\calM,\quad
\dot\Sigma_{\calC}=-\beta_0 j_0-\beta_1 j_1-\beta_z j_{\calC},\quad
\dot\Sigma_\calM=-\beta_z j_\calM-\beta_r j_r,
\end{equation}
with $j_{0,1,r}$ the currents from $B_{0,1,r}$ into their qubits.
Integrating over the run time $\tau$ and averaging over inputs [paper
Eqs.~(28)--(29)],
\begin{equation}
\langle\Sigma\rangle=\sum_{x\in\{0,1\}}\int p(x)\,e(\beta_1\mid x)\,
\Sigma(\beta_1)\,\dd\beta_1,\qquad
\Sigma(\beta_1)=\int_0^\tau\dot\Sigma(\beta_1)\,\dd t.
\label{eq:diss}
\end{equation}

\begin{keybox}
\textbf{The central trade--off} [paper Fig.~3C]. Lowering the average
error $\langle\xi\rangle$ requires increasing the gain, i.e.\ raising
$\epsilon_1$ (Eq.~\eqref{eq:gain}); but the heat currents --- and hence
$\langle\Sigma\rangle$ --- grow with the qubit energies. Thus
$\langle\xi\rangle$ and $\langle\Sigma\rangle$ are monotonically
opposed: \emph{accuracy costs dissipation}. This is the thermodynamic
price of computation for the neuron, and the baseline that the
Floquet--buffer extension seeks to improve.
\end{keybox}

\begin{keybox}
\textbf{End of Part III.} Full machine Hamiltonian:
$H=H_0+H_{\mathrm{int}}+H_\calM$ (collector free $+$ resonant swap $+$
modulator), plus four local reset baths. The collector builds the
weighted sum $\beta_v$; the modulator saturates it into the sigmoid
\eqref{eq:sigmoid-final}; encoding \eqref{eq:encoding} reads the bit; and
accuracy trades against dissipation. Part IV lifts this to $n$ inputs and
proves the perceptron equivalence.
\end{keybox}

% #####################################################################
\part{The general neuron: perceptron equivalence and networks}
\label{part:general}
% #####################################################################

\section{The $n$--input collector}
\label{sec:n-collector}
Everything generalizes. The collector now has $n+2$ qubits: a reference
$C_0$ (bath $B_0$, fixed $\beta_0$), inputs $C_1,\dots,C_n$ (baths
$B_1,\dots,B_n$ at $\beta_1,\dots,\beta_n$ encoding the $n$ input bits),
and the output $C_z$ (finite reservoir $B_z$). The modulator is
\emph{exactly the same} single qubit as before. The only new ingredient
is that a multi--qubit machine hosts \emph{many} virtual qubits, and we
must say which one we use.

\subsection{Choosing a virtual qubit: the vector $h$}
Introduce a binary vector $h=(h_0,h_1,\dots,h_n)$, $h_i\in\{0,1\}$,
recording whether qubit $i$ contributes its ground ($h_i=0$) or excited
($h_i=1$) state. The chosen virtual qubit has levels [paper
Eqs.~(30)--(31)]
\begin{align}
\ket{0}_v &:= \ket{h_0}_{C_0}\ket{h_1}_{C_1}\cdots\ket{h_n}_{C_n},\\
\ket{1}_v &:= \ket{h_0\!\oplus\!1}_{C_0}\ket{h_1\!\oplus\!1}_{C_1}
\cdots\ket{h_n\!\oplus\!1}_{C_n},
\end{align}
where $\oplus$ is addition mod $2$: $\ket{1}_v$ is the bit--flip of
$\ket{0}_v$ on \emph{every} qubit. Its energy gap is [paper Eq.~(32)]
\begin{equation}
\epsilon_v = \sum_{i=0}^{n}(-1)^{h_i\oplus1}\epsilon_i .
\label{eq:eps-v-gen}
\end{equation}

\subsection{The interaction and the multi--input virtual temperature}
As before, choose $C_z$ resonant with this virtual qubit,
$\epsilon_z=\epsilon_v$, and couple them by the swap [paper, after
Eq.~(32)]
\begin{equation}
H_{\mathrm{int}} = g\big(\ket{0}\bra{1}_v\otimes\ket{1}\bra{0}_{C_z}
+ \mathrm{h.c.}\big),
\label{eq:Hint-gen}
\end{equation}
which is the direct generalization of \eqref{eq:Hint-swap}. The
excited--state population of the virtual qubit is the product of the
per--qubit occupations [paper Eq.~(33)]
\begin{equation}
g_v = g_0(h_0)\,g_1(h_1)\cdots g_n(h_n),\qquad
g_i(0)=(1+\ee^{-\beta_i\epsilon_i})^{-1},\ g_i(1)=1-g_i(0),
\end{equation}
and $\ee^{-\beta_v\epsilon_v}=g_v/(1-g_v)$ gives, after taking the log of
the product [paper Eq.~(34)],
\begin{equation}
\boxed{\;
\beta_v = \frac{1}{\epsilon_z}\sum_{i=0}^{n}(-1)^{h_i}\beta_i\epsilon_i .
\;}
\label{eq:betav-gen}
\end{equation}
Again a \emph{linear combination of the input inverse temperatures}, with
weights $(-1)^{h_i}\epsilon_i/\epsilon_z$ set by the gaps and the vector
$h$. The steady--state output is the same nonlinear map as in Part III,
[paper Eqs.~(35)--(36)]
\begin{equation}
\beta_z^\infty=\frac{1}{\epsilon_z}\ln\!\big[Q(\beta_v)^{-1}-1\big]
= f(\beta_v)+O(\epsilon_z),
\end{equation}
with $Q$ from \eqref{eq:transfer}. The whole machine still computes
``weighted sum, then sigmoid.''

\section{Exact correspondence with the perceptron}
\label{sec:perceptron}
A perceptron with inputs $x=(x_0,\dots,x_n)$ (with $x_0=1$ by convention)
and weights $w=(w_0,\dots,w_n)$ outputs [paper Eq.~(37)]
\begin{equation}
z=f(y),\qquad y=\sum_{i=0}^{n}w_i x_i .
\end{equation}
The two--step structure matches the neuron line for line:
\begin{center}
\begin{tabular}{@{}ll@{}}
\toprule
Perceptron & Thermodynamic neuron \\
\midrule
inputs $x_i$ & input inverse temperatures $\beta_i$ \\
weighted sum $y=\sum_i w_i x_i$ & virtual temperature $\beta_v$,
Eq.~\eqref{eq:betav-gen} \\
bias $w_0$ & reference term $(-1)^{h_0}\beta_0\epsilon_0/\epsilon_z$ \\
weights $w_i$ & $(-1)^{h_i}\epsilon_i/\epsilon_z$ \\
activation $f$ (sigmoid) & modulator map, small--$\epsilon_z$
Eq.~\eqref{eq:sigmoid-final} \\
output $z$ & output temperature $\beta_z^\infty$ \\
\bottomrule
\end{tabular}
\end{center}
Concretely, with an overall energy scale $\alpha$ (defined below), the
pre--activation is [paper Eq.~(44)]
\begin{equation}
\beta_z^\infty=f(x)+O(\epsilon_z),\qquad
x=\alpha\Big(w_0+\sum_{i=1}^{n}w_i\beta_i\Big),
\label{eq:perceptron-out}
\end{equation}
\emph{exactly} the sigmoid perceptron. Hence a single thermodynamic
neuron implements precisely the linearly--separable functions, and (by
choosing other thermalization models) a strict superset.

\section{Designing a neuron: the algorithm}
\label{sec:algorithm}
Because the neuron \emph{is} a perceptron, we design it by training a
perceptron and reading off the hardware. Given a target $n$--input
Boolean function $R(x)$ [paper, Algorithm 1]:
\begin{enumerate}[leftmargin=*]
\item Build the training set $D=\{(x^{(i)},R(x^{(i)}))\}_{i=1}^{2^n}$.
\item Train a linear classifier (sigmoid perceptron) to obtain weights
$w=(w_0,\dots,w_n)$.
\item Set the interaction vector from the signs of the weights,
$h_k = 0$ if $w_k\ge0$, else $h_k=1$ [paper Eq.~(39)].
\item Set the qubit energies [paper Eq.~(40)]
\begin{equation}
\epsilon_k=\begin{cases}
\alpha\,\big|\epsilon_z-\sum_{k=1}^n w_k\big| & k=0,\\[2pt]
\alpha\,|w_k| & \text{otherwise.}
\end{cases}
\end{equation}
\item Set the reference (bias) temperature [paper Eq.~(41)]
\begin{equation}
\beta_0=\frac{|w_0|}{\big|\epsilon_z-\sum_{k=1}^n w_k\big|}.
\end{equation}
\end{enumerate}
Substituting into \eqref{eq:betav-gen} reproduces
$\beta_v=\tfrac{\alpha}{\epsilon_z}\big(w_0+\sum_k w_k\beta_k\big)$
[paper Eqs.~(42)--(43)], i.e.\ the trained hyperplane. The scale
$\alpha>0$ is the ``quality knob'': it rescales all gaps, sharpening the
threshold (better classification) at the price of proportionally more
dissipation --- the same accuracy/dissipation trade--off of
Sec.~\ref{sec:tradeoff}, now global. For the NOT gate one finds
$\alpha=\epsilon_1$.

\begin{implbox}[design + simulate a neuron]
Train any logistic/perceptron classifier on the truth table (e.g.\
\texttt{sklearn}); take \texttt{w}; map to \texttt{h, eps\_k, beta0} by
the boxed formulas; build $H_0+H_{\mathrm{int}}$ (with the generalized
swap \eqref{eq:Hint-gen}) plus reset baths in QuTiP; integrate the
two--tier dynamics of Sec.~\ref{sec:dynamics}; read $\beta_z^\infty$ and
decode with \eqref{eq:encoding}. Sweeping the inputs reproduces the
response surfaces of [paper Figs.~6--7].
\end{implbox}

\section{Worked examples}
\label{sec:examples}
\begin{example}[NOR, $n=2$, paper Eqs.~(45)--(46)]
The separating hyperplane $x_1+x_2=\tfrac12$ gives weights
$w=(1,-2,-2)$. Then
\begin{equation}
h=(0,1,1),\qquad \epsilon=\alpha(\epsilon_z+4,\,2,\,2),\qquad
\beta_0=(\epsilon_z+4)^{-1},
\end{equation}
and the virtual temperature becomes
$\beta_v=\alpha(1-2\beta_1-2\beta_2)$. The response
$\beta_z^\infty(\beta_1,\beta_2)$ matches the NOR truth table. Because
NOR is functionally complete, networks of these neurons are universal.
\end{example}

\begin{example}[3--MAJORITY, $n=3$, paper Eqs.~(47)--(48)]
Training gives $w=(-4,3,3,3)$, hence
\begin{equation}
h=(1,0,0,0),\qquad \epsilon=\alpha(\epsilon_z+12,\,3,\,3,\,3),\qquad
\beta_0=(\epsilon_z+12)^{-1},
\end{equation}
and $\beta_v=\alpha(4-3\beta_1-3\beta_2-3\beta_3)$. The response
reproduces the majority function.
\end{example}

\section{Networks: beyond linear separability (XOR)}
\label{sec:networks}
A single neuron can only realize linearly--separable functions (the
perceptron limitation). XOR is not separable, so it needs a
\emph{network}: write $\text{XOR}=\text{AND}(\text{NAND},\text{OR})$ and
wire three neurons accordingly [paper, Fig.~8].

\paragraph{The wiring problem and its fix.}
Feeding one neuron's \emph{finite} output bath $B_z$ into the next
neuron's input would let unwanted (possibly back--flowing) currents
corrupt the operation, and Eq.~\eqref{eq:transfer} would no longer hold.
Two remedies:
\begin{itemize}[leftmargin=*]
\item \emph{Measurement gateway:} an external agent reads $\beta_z$ of
the first neuron and couples the second neuron's input qubit to an
\emph{infinite} bath at $\beta_{\mathrm{hot}}$ or $\beta_{\mathrm{cold}}$
accordingly. Simple, but sacrifices autonomy.
\item \emph{Clock gateway} [paper, Supp.~C]: a stopwatch alternately
opens/closes the coupling between $B_z$ and the next input qubit. Define
$t_{\mathrm{steady}}$ (time for neuron~1 to reach $\beta_z^\infty$ with
the link \emph{off}) and $t_{\mathrm{const}}$ (interval over which $B_z$
stays effectively constant when used as neuron~2's input). Turning the
link on for $t_{\mathrm{const}}$ after each $t_{\mathrm{steady}}$ makes
$B_z$ behave as an infinite bath to neuron~2, with no spurious
back--currents. Realized with an autonomous thermal clock, the network
stays fully autonomous.
\end{itemize}
This clock--gated readout is exactly the mechanism the Floquet--buffer
extension internalizes: the $J_0$--zero gating of the output buffer plays
the role of $t_{\mathrm{const}}/t_{\mathrm{steady}}$, replacing the
external stopwatch by the drive itself.

\begin{keybox}
\textbf{End of Part IV.} The $n$--input neuron computes
$\beta_z^\infty\approx f\!\big(\alpha(w_0+\sum_i w_i\beta_i)\big)$, a
sigmoid perceptron; its hardware ($h,\epsilon_k,\beta_0$) is read off a
trained classifier; networks (clock--gated) reach any Boolean function.
This is the complete autonomous baseline; the Floquet buffer aims to
improve its accuracy/dissipation frontier and to make composition
native.
\end{keybox}

\appendix
\part*{Appendices}
\addcontentsline{toc}{part}{Appendices}

\section{Notation and symbol table}
\label{app:notation}
\begin{center}
\begin{tabularx}{\textwidth}{@{}llX@{}}
\toprule
Symbol & Meaning & First appears \\
\midrule
$\epsilon_i$ & energy gap of qubit $C_i$ & \eqref{eq:H-single} \\
$\beta_i$ & inverse temperature of bath $B_i$ & \eqref{eq:gibbs} \\
$g(\beta\epsilon)$ & Fermi--Dirac / sigmoid occupation & \eqref{eq:fermi} \\
$\tau(\beta)$ & Gibbs state & \eqref{eq:gibbs} \\
$\calL^{(k)}$ & reset dissipator for bath $k$ & \eqref{eq:reset} \\
$\gamma_k,\mu,\mu'$ & thermalization rates & \eqref{eq:reset}, \eqref{eq:jC} \\
$\beta_v,\epsilon_v$ & virtual temperature and gap & \eqref{eq:eps-v}--\eqref{eq:betav} \\
$\chi$ & collector coupling strength & \eqref{eq:Hint-kets} \\
$j_k$ & heat current into bath $k$ & \eqref{eq:current-def} \\
$\dot\Sigma$ & entropy production rate & \eqref{eq:entropy-prod} \\
\bottomrule
\end{tabularx}
\end{center}

\section{List of approximations}
\label{app:approx}
\begin{center}
\begin{tabularx}{\textwidth}{@{}llX@{}}
\toprule
Tag & Name & Regime of validity \\
\midrule
(A1) & Born / weak coupling & $\gamma\ll\epsilon$, factorized state \\
(A2) & Markov / short memory & bath correlation time $\tau_c\ll1/\gamma$ \\
(A3) & secular / RWA & $\gamma\ll\epsilon$ (resolved gap) \\
(A4) & reset thermalization & steady--state populations only \\
(A5) & energy--preserving $H_{\mathrm{int}}$ & $\chi\ll\epsilon$ \\
(A6) & local master equation & $\chi\ll\epsilon$ (else global) \\
\bottomrule
\end{tabularx}
\end{center}

\section{Full Born--Markov--secular derivation}
\label{app:BMS}
We fill in the steps compressed in Sec.~\ref{sec:BMS}. Start from
\eqref{eq:exact-integro} and trace over the bath. Writing
$H_{SB}=\sum_\alpha A_\alpha\otimes B_\alpha$ (here $A=\sigma_x$, one
term), the second--order generator is
\begin{equation}
\dot{\tilde\rho}(t) = -\int_0^t\dd s\,\Tr_B\!\big[\tilde H_{SB}(t),
[\tilde H_{SB}(s),\tilde\rho(s)\otimes\rho_B]\big].
\end{equation}
Expanding the double commutator and using
$\Tr_B[B_\alpha\rho_B]=0$ leaves products of the \emph{bath correlation
functions}
\begin{equation}
C_{\alpha\beta}(t-s)=\Tr_B\!\big[\tilde B_\alpha(t)\tilde B_\beta(s)
\rho_B\big],
\end{equation}
which depend only on $\tau=t-s$ because $\rho_B$ is stationary. Decompose
the system coupling operator into eigenoperators of $H_S$,
\begin{equation}
A(\omega)=\sum_{\varepsilon'-\varepsilon=\omega}
\Pi(\varepsilon)\,A\,\Pi(\varepsilon'),\qquad
\tilde A(\omega,t)=\ee^{-\ii\omega t}A(\omega),
\end{equation}
with $\Pi$ the projectors onto the eigenspaces of $H_S$; for the qubit
$A(\pm\epsilon)=\sigma_\mp$. Making the Markov replacement (A2)
$\tilde\rho(s)\to\tilde\rho(t)$ and extending the integral to infinity,
one gets the Redfield equation with the half--Fourier rates
\begin{equation}
\Gamma(\omega)=\int_0^\infty\dd\tau\,\ee^{\ii\omega\tau}C(\tau)
= \tfrac12\gamma(\omega)+\ii S(\omega),
\end{equation}
where $\gamma(\omega)$ is the full Fourier transform
\eqref{eq:rate-fourier} (a real, positive relaxation rate) and
$S(\omega)$ its imaginary part (the Lamb shift). The secular
approximation (A3) drops all cross terms with $\omega\neq\omega'$
(they rotate at $\ee^{\ii(\omega'-\omega)t}$), yielding the diagonal,
completely positive GKSL form
\begin{equation}
\dot\rho=-\ii[H_S+H_{LS},\rho]
+\sum_{\omega}\gamma(\omega)\Big(A(\omega)\rho A^\dagger(\omega)
-\tfrac12\{A^\dagger(\omega)A(\omega),\rho\}\Big),
\end{equation}
with $H_{LS}=\sum_\omega S(\omega)A^\dagger(\omega)A(\omega)$. For the
qubit this is exactly \eqref{eq:GKSL-qubit} with
$\gamma_\downarrow=\gamma(\epsilon)$, $\gamma_\uparrow=\gamma(-\epsilon)$,
and the KMS relation $\gamma(-\omega)=\ee^{-\beta\omega}\gamma(\omega)$
gives \eqref{eq:db-rates}. The Lamb shift $H_{LS}\propto\sigma_+\sigma_-$
merely renormalizes $\epsilon$ and is usually absorbed into it. See
Ref.~\cite{Breuer2002}, Ch.~3, for the full treatment.

\section{The reset model as a collisional process}
\label{app:reset}
The reset dissipator \eqref{eq:reset} arises from a repeated--interaction
(collisional) model. In each interval $\dd t$, qubit $k$ either (prob.\
$\gamma_k\dd t$) is swapped for a fresh bath qubit in state
$\tau(\beta_k)$, or (prob.\ $1-\gamma_k\dd t$) is left untouched:
\begin{equation}
\rho(t+\dd t) = (1-\gamma_k\dd t)\,\rho(t)
+ \gamma_k\dd t\,\big(\Tr_k[\rho(t)]\otimes\tau(\beta_k)\big).
\end{equation}
Rearranging and taking $\dd t\to0$ gives exactly
$\dot\rho=\gamma_k(\Tr_k[\rho]\otimes\tau(\beta_k)-\rho)$. The map is
manifestly completely positive and trace preserving, and has
$\tau(\beta_k)$ as its unique fixed point on qubit $k$, so it satisfies
local detailed balance \eqref{eq:ldb} by construction. This is why the
reset model can stand in for the microscopically derived dissipator
whenever only steady--state populations and their thermodynamics matter
(approximation (A4)).

\section{Numerical implementation in QuTiP}
\label{app:numerics}
A minimal, self--contained recipe that reproduces the NOT gate.

\paragraph{Operators.} With ordering $(C_0C_1C_z)$:
\begin{verbatim}
import numpy as np, qutip as qt
I2 = qt.qeye(2); sm = qt.destroy(2); sp = sm.dag(); n = sp*sm
def op(o,site): L=[I2,I2,I2]; L[site]=o; return qt.tensor(L)
n0,n1,nz = op(n,0),op(n,1),op(n,2)
H0   = eps0*n0 + eps1*n1 + eps_z*nz              # free Hamiltonian, Eq.(A? / HS)
k100 = qt.tensor(qt.basis(2,1),qt.basis(2,0),qt.basis(2,0))
k011 = qt.tensor(qt.basis(2,0),qt.basis(2,1),qt.basis(2,1))
Hint = chi*(k100*k011.dag() + k011*k100.dag())   # resonant swap, Eq.(Hint)
H    = H0 + Hint
assert (qt.commutator(H0,Hint)).norm() < 1e-12    # resonance sanity check
\end{verbatim}

\paragraph{Reset dissipators.} Build collapse operators reproducing the
reset fixed point for each bath (rate $\gamma_i$, occupation
$g_i=g(\beta_i\epsilon_i)$):
\begin{verbatim}
def bath(site, gamma, g):     # detailed-balanced up/down channels
    return [np.sqrt(gamma*(1-g))*op(sm,site),   # emission  (down)
            np.sqrt(gamma*g)    *op(sp,site)]   # absorption(up)
g0 = 1/(1+np.exp(beta0*eps0)); g1 = 1/(1+np.exp(beta1*eps1))
gz = 1/(1+np.exp(betaz*eps_z))
c_ops = bath(0,gamma,g0) + bath(1,gamma,g1) + bath(2,mu,gz)
rho_ss = qt.steadystate(H, c_ops)                # collector-only NESS
p_z = qt.expect(nz, rho_ss)                       # excited pop of C_z
beta_z = np.log((1-p_z)/p_z)/eps_z                # -> approaches beta_v
\end{verbatim}

\paragraph{Two--tier dynamics (the full gate).} Because
$\gamma\gg\mu,\mu'$ (Sec.~\ref{sec:dynamics}), do not integrate the stiff
full model directly. Instead, for a grid of $\beta_z$ values compute the
fast steady--state currents
$j_{\calC}(\beta_z)=\mu\epsilon_z[\gz(\beta_z)-\gz(\beta_v)]$ and
$j_\calM(\beta_z)=\mu'\epsilon_z[\gz(\beta_z)-\gz(\beta_r)]$, then
integrate the scalar slow ODE \eqref{eq:calorimetric}
$\dot\beta_z=-(\beta_z^2/\calC)(j_{\calC}+j_\calM)$ with a standard
(non--stiff) solver; its fixed point is \eqref{eq:transfer}. Sweeping
$\beta_1\in\{\beta_{\mathrm{hot}},\beta_{\mathrm{cold}}\}$ and decoding
with \eqref{eq:encoding} yields the truth table; sweeping $\beta_1$
continuously yields the transfer curve [paper Fig.~3B].

\paragraph{Currents, entropy production, error.} Heat currents follow
from \eqref{eq:current-def}, $j_k=\Tr[H\,\calL^{(k)}\rho]$; entropy
production from \eqref{eq:entropy-prod},
$\dot\Sigma=-\sum_k\beta_k j_k$ in the NESS. The readout error
\eqref{eq:error} is obtained by convolving the transfer curve with the
Gaussian $\mathcal N(\beta_z^\infty,C)$ and integrating against the guard
bands. Plotting $\langle\Sigma\rangle$ vs $\langle\xi\rangle$ over a
sweep of $\epsilon_1$ reproduces the trade--off of [paper Fig.~3C].

\begin{implbox}[validation checklist]
(i) $\|[H_0,H_{\mathrm{int}}]\|=0$ (resonance). (ii) Collector--only
steady state gives $\beta_z\to\beta_v$ within tolerance. (iii) First law
per bath: $\sum_k j_k=0$ in the NESS (energy conservation). (iv)
$\dot\Sigma\ge0$ at every input. (v) Transfer curve from the ODE matches
the closed form \eqref{eq:transfer}. Only after all five pass should any
physics be trusted.
\end{implbox}

\begin{thebibliography}{9}
\bibitem{LipkaBartosik2024}
P.~Lipka--Bartosik, M.~Perarnau--Llobet, and N.~Brunner,
``Thermodynamic computing via autonomous quantum thermal machines,''
\emph{Sci.\ Adv.} \textbf{10}, eadm8792 (2024); arXiv:2308.15905.
\bibitem{BrunnerVirtual2012}
N.~Brunner, N.~Linden, S.~Popescu, and P.~Skrzypczyk,
``Virtual qubits, virtual temperatures, and the foundations of
thermodynamics,'' \emph{Phys.\ Rev.\ E} \textbf{85}, 051117 (2012).
\bibitem{Linden2010}
N.~Linden, S.~Popescu, and P.~Skrzypczyk,
``How small can thermal machines be? The smallest possible
refrigerator,'' \emph{Phys.\ Rev.\ Lett.} \textbf{105}, 130401 (2010).
\bibitem{Hofer2017}
P.~P.~Hofer \emph{et al.},
``Markovian master equations for quantum thermal machines: local versus
global approach,'' \emph{New J.\ Phys.} \textbf{19}, 123037 (2017).
\bibitem{Breuer2002}
H.-P.~Breuer and F.~Petruccione,
\emph{The Theory of Open Quantum Systems}, Oxford University Press
(2002).
\end{thebibliography}

\end{document}
